\chapter{高等数学：积分不等式}

\section{积分不等式的基本工具}

\begin{problemblock}
\textbf{例1.} 设 \(f(x)\) 在 \([0,1]\) 上可导，当 \(0\le x\le1\) 时，
\[
f'(x)+f^2(x)\ge0,\qquad f(0)>0,
\]
则（\quad）

A. \(\displaystyle\int_0^1f(x)\,dx\le\ln\frac{f(1)}{f(0)}\)

B. \(\displaystyle\int_0^1f(x)\,dx\ge\ln\frac{f(0)}{f(1)}\)

C. \(\displaystyle\int_0^1f(x)\,dx\le\ln f(1)\)

D. \(\displaystyle\int_0^1f(x)\,dx\ge\ln f(0)\)
\end{problemblock}
\begin{solutionblock}
\analysis{题设中出现 \(f'\) 与 \(f^2\)，应想到除以 \(f\) 或 \(f^2\)。先要证明 \(f\) 在 \([0,1]\) 上不变号，否则对数无意义。}
因为 \(f(0)>0\)。若 \(f\) 在 \([0,1]\) 上第一次取到 \(0\) 的点为 \(c>0\)，则在 \([0,c)\) 上 \(f>0\)。由
\[
f'(x)+f^2(x)\ge0
\]
得
\[
\left(\frac1{f(x)}\right)'=-\frac{f'(x)}{f^2(x)}\le1.
\]
积分可知 \(1/f(x)\) 在 \(x\to c^-\) 时不可能趋于 \(+\infty\)，矛盾。因此
\[
f(x)>0,\qquad 0\le x\le1.
\]
在正值条件下，将原不等式除以 \(f(x)\)，得
\[
\frac{f'(x)}{f(x)}+f(x)\ge0,
\]
即
\[
f(x)\ge-\frac{f'(x)}{f(x)}.
\]
两边在 \([0,1]\) 上积分：
\[
\int_0^1 f(x)\,dx
\ge-\int_0^1\frac{f'(x)}{f(x)}\,dx
=-\ln f(x)\bigg|_0^1
=\ln\frac{f(0)}{f(1)}.
\]
\answer{B}
\examnote{遇到 \(f'+f^2\ge0\) 且 \(f(0)>0\)，先证明 \(f>0\)，再用 \((\ln f)'=f'/f\)。}
\end{solutionblock}

\begin{problemblock}
\textbf{例2.} 设函数 \(f(x)\) 有一阶连续导数，且
\[
0<f'(x)\le1,\qquad f(0)=0.
\]
证明：
\[
\left[\int_0^1f(x)\,dx\right]^2
\ge
\int_0^1[f(x)]^3\,dx.
\]
\end{problemblock}
\begin{solutionblock}
\analysis{关键是由 \(f'\le1\) 推出 \(f^2(x)\le2\int_0^x f(t)\,dt\)，再构造“上限函数”证明单调。}
因 \(f'(x)>0\) 且 \(f(0)=0\)，所以
\[
f(x)>0\quad (x>0).
\]
记
\[
F(x)=\int_0^x f(t)\,dt.
\]
由 \(0<f'\le1\)，有
\[
f^2(x)=2\int_0^x f(t)f'(t)\,dt
\le2\int_0^x f(t)\,dt=2F(x).
\]
再令
\[
H(x)=F^2(x)-\int_0^x f^3(t)\,dt.
\]
则
\[
H'(x)=2F(x)f(x)-f^3(x)
=f(x)\bigl[2F(x)-f^2(x)\bigr]\ge0.
\]
又 \(H(0)=0\)，故 \(H(1)\ge0\)，即
\[
\left[\int_0^1f(x)\,dx\right]^2
\ge
\int_0^1f^3(x)\,dx.
\]
\answer{不等式成立。}
\examnote{证明积分不等式时，常把目标差写成 \(H(1)-H(0)\)，再证明 \(H'(x)\ge0\)。}
\end{solutionblock}

\begin{problemblock}
\textbf{例3.} 设 \(f(x)\) 是 \([0,1]\) 上单调增加的连续函数，则（\quad）

A. \(\displaystyle \int_0^{\int_0^1e^{-t^2}\,dt}f(x)\,dx
\ge\int_0^1f(x)e^{-x^2}\,dx\)

B. \(\displaystyle \int_0^{\int_0^1e^{-t^2}\,dt}f(x)\,dx
\le\int_0^1f(x)e^{-x^2}\,dx\)

C. \(\displaystyle \int_0^{\int_0^1e^{-t^2}\,dt}f(x)\,dx
\ge\int_0^1f(x)\,dx\)

D. \(\displaystyle \int_0^{\int_0^1e^{-t^2}\,dt}f(x)\,dx
\le\int_0^1f(x)\,dx\)
\end{problemblock}
\begin{solutionblock}
\analysis{权函数 \(e^{-x^2}\) 单调递减，且它的总质量等于左边积分区间长度。对递增函数，质量越偏左，积分越小。}
记
\[
\alpha=\int_0^1e^{-t^2}\,dt,\qquad 0<\alpha<1.
\]
需要比较
\[
L=\int_0^\alpha f(x)\,dx,\qquad
R=\int_0^1f(x)e^{-x^2}\,dx.
\]
有
\[
R-L
=\int_0^\alpha f(x)(e^{-x^2}-1)\,dx
+\int_\alpha^1 f(x)e^{-x^2}\,dx.
\]
又
\[
\int_0^\alpha(1-e^{-x^2})\,dx
=\int_\alpha^1e^{-x^2}\,dx
\]
由 \(\alpha=\int_0^1e^{-x^2}dx\) 得到。由于 \(f\) 单调增加，
\[
f(x)\le f(\alpha)\quad(0\le x\le\alpha),
\qquad
f(x)\ge f(\alpha)\quad(\alpha\le x\le1).
\]
因此
\[
R-L
\ge -f(\alpha)\int_0^\alpha(1-e^{-x^2})\,dx
+f(\alpha)\int_\alpha^1e^{-x^2}\,dx=0.
\]
所以 \(L\le R\)。
\answer{B}
\examnote{这是“等质量重排”思想：递增函数配递减权时，权重向右分散会使积分变大。}
\end{solutionblock}

\begin{problemblock}
\textbf{例4.} 设
\[
I_n=\int_0^{\frac\pi4}\tan^n x\,dx,\qquad n\in\mathbb{N}^+,\ n\ge2.
\]
证明：
\[
\frac1{2(n+1)}<I_n<\frac1{2(n-1)}.
\]
\end{problemblock}
\begin{solutionblock}
\analysis{令 \(t=\tan x\) 后，积分区间变成 \([0,1]\)，被积函数变为 \(t^n/(1+t^2)\)，再用简单夹逼。}
令 \(t=\tan x\)，则 \(x\in[0,\pi/4]\) 对应 \(t\in[0,1]\)，且
\[
dx=\frac{dt}{1+t^2}.
\]
于是
\[
I_n=\int_0^1\frac{t^n}{1+t^2}\,dt.
\]
当 \(0<t<1\) 时，
\[
\frac12<\frac1{1+t^2}<\frac1{2t^2}.
\]
乘以 \(t^n\) 得
\[
\frac12t^n<\frac{t^n}{1+t^2}<\frac12t^{n-2}.
\]
两边积分：
\[
\frac12\int_0^1t^n\,dt<I_n<\frac12\int_0^1t^{n-2}\,dt.
\]
因此
\[
\frac1{2(n+1)}<I_n<\frac1{2(n-1)}.
\]
\answer{不等式成立。}
\examnote{\(\tan x\) 在 \(0\) 到 \(\pi/4\) 上常令 \(t=\tan x\)，把三角积分转成 \([0,1]\) 上的有理积分。}
\end{solutionblock}

\begin{problemblock}
\textbf{例5.} （切比雪夫不等式）设 \(f(x),g(x)\) 在 \([a,b]\) 上连续，且 \(f(x),g(x)\) 在 \([a,b]\) 上单调性相同；\(p(x)\ge0\) 且在 \([a,b]\) 上连续。证明：
\[
\int_a^bp(x)f(x)\,dx\int_a^bp(x)g(x)\,dx
\le
\int_a^bp(x)\,dx\int_a^bp(x)f(x)g(x)\,dx.
\]
\end{problemblock}
\begin{solutionblock}
\analysis{切比雪夫不等式的核心是同单调函数满足 \((f(x)-f(y))(g(x)-g(y))\ge0\)。用双重积分最直接。}
若 \(p\equiv0\)，结论显然。下设 \(p\) 不恒为零。考虑
\[
\Delta=
\int_a^bp(x)\,dx\int_a^bp(x)f(x)g(x)\,dx
-\int_a^bp(x)f(x)\,dx\int_a^bp(x)g(x)\,dx.
\]
把乘积写成二重积分：
\[
\begin{aligned}
2\Delta
&=\int_a^b\int_a^b p(x)p(y)
\bigl[f(x)-f(y)\bigr]\bigl[g(x)-g(y)\bigr]\,dx\,dy.
\end{aligned}
\]
因为 \(p(x)p(y)\ge0\)，且 \(f,g\) 单调性相同，所以
\[
\bigl[f(x)-f(y)\bigr]\bigl[g(x)-g(y)\bigr]\ge0.
\]
故 \(2\Delta\ge0\)，即 \(\Delta\ge0\)，所证不等式成立。
\answer{不等式成立。}
\examnote{同单调用切比雪夫，反单调用反向切比雪夫；二重积分恒等式是最稳的证明模板。}
\end{solutionblock}

\section{比较法与导数约束}

\begin{problemblock}
\textbf{例6.} 证明：
\[
\int_0^{\frac\pi2}\frac{\cos x}{1+\ln(1+x)}\,dx
\ge
\int_0^{\frac\pi2}\frac{\sin x}{1+\ln(1+x)}\,dx.
\]
\end{problemblock}
\begin{solutionblock}
\analysis{把两边相减后，利用 \(x\) 与 \(\pi/2-x\) 配对；分母对应的权函数是递减的。}
令
\[
h(x)=\frac1{1+\ln(1+x)}.
\]
则 \(h(x)\) 在 \([0,\pi/2]\) 上单调递减。两边差为
\[
D=\int_0^{\frac\pi2}(\cos x-\sin x)h(x)\,dx.
\]
配对 \(x\) 与 \(\frac\pi2-x\)，得
\[
D=\int_0^{\frac\pi4}
(\cos x-\sin x)\left[h(x)-h\left(\frac\pi2-x\right)\right]dx.
\]
在 \(0\le x\le\pi/4\) 上，
\[
\cos x-\sin x\ge0,
\qquad
x\le\frac\pi2-x.
\]
又 \(h\) 单调递减，所以
\[
h(x)-h\left(\frac\pi2-x\right)\ge0.
\]
于是 \(D\ge0\)，原不等式成立。
\answer{不等式成立。}
\examnote{三角积分比较题常把 \(\sin x\) 与 \(\cos x\) 通过 \(x\mapsto\pi/2-x\) 配对。}
\end{solutionblock}

\begin{problemblock}
\textbf{例7.} 已知
\[
I_1=\int_0^{\frac\pi2}\frac{\sin x}{1+x^2}\,dx,\qquad
I_2=\int_0^{\frac\pi2}\frac{\cos x}{1+x^2}\,dx,
\]
\[
J_1=\int_0^1 e^x\sqrt{1-x}\,dx,\qquad
J_2=\int_0^1 e^x\frac{x}{2\sqrt{1-x}}\,dx,
\]
则（\quad）

A. \(I_1>I_2,\ J_1>J_2\)\qquad
B. \(I_1>I_2,\ J_1<J_2\)

C. \(I_1<I_2,\ J_1>J_2\)\qquad
D. \(I_1<I_2,\ J_1<J_2\)
\end{problemblock}
\begin{solutionblock}
\analysis{前一组用 \(\sin,\cos\) 配对；后一组用分部积分找出 \(J_1,J_2\) 的关系。}
先比较 \(I_1,I_2\)。令 \(h(x)=1/(1+x^2)\)，它在 \([0,\pi/2]\) 上递减。于是
\[
I_2-I_1
=\int_0^{\frac\pi2}(\cos x-\sin x)h(x)\,dx
\]
配对后为
\[
\int_0^{\frac\pi4}(\cos x-\sin x)
\left[h(x)-h\left(\frac\pi2-x\right)\right]dx>0.
\]
故
\[
I_1<I_2.
\]

再比较 \(J_1,J_2\)。记
\[
K=\int_0^1\frac{e^x}{2\sqrt{1-x}}\,dx.
\]
对 \(J_1\) 分部积分，取 \(u=\sqrt{1-x}\)，\(dv=e^x\,dx\)，得
\[
J_1=\left[e^x\sqrt{1-x}\right]_0^1+K=-1+K,
\]
所以 \(K=J_1+1\)。又
\[
K-J_2=\int_0^1 e^x\frac{1-x}{2\sqrt{1-x}}\,dx
=\frac12J_1.
\]
因此
\[
J_2=K-\frac12J_1=1+\frac12J_1.
\]
由于
\[
J_1<\int_0^1 e^x\,dx=e-1<2,
\]
所以
\[
J_2=1+\frac12J_1>J_1.
\]
故 \(J_1<J_2\)。
\answer{D}
\examnote{比较两个看似相近的积分，若直接点态比较困难，分部积分常能制造出线性关系。}
\end{solutionblock}

\begin{problemblock}
\textbf{例8.} 设函数 \(f(x)\) 在 \([a,b]\) 上不恒为零，且
\[
f(a)=f(b)=0,\qquad f'(x)\ \text{在}\ [a,b]\ \text{上连续}.
\]
试证：至少存在一点 \(\xi\in[a,b]\)，使得
\[
|f'(\xi)|\ge
\frac1{(b-a)^2}\int_a^b f(x)\,dx.
\]
\end{problemblock}
\begin{solutionblock}
\analysis{可以证明更强的不等式。令 \(M=\max|f'|\)，利用端点为零估计 \(|f(x)|\) 的高度。}
由 \(f'\) 连续，存在
\[
M=\max_{x\in[a,b]}|f'(x)|.
\]
对任意 \(x\in[a,b]\)，由 \(f(a)=0\) 得
\[
|f(x)|\le M(x-a),
\]
由 \(f(b)=0\) 得
\[
|f(x)|\le M(b-x).
\]
因此
\[
|f(x)|\le M\min\{x-a,b-x\}.
\]
积分得
\[
\left|\int_a^b f(x)\,dx\right|
\le\int_a^b|f(x)|\,dx
\le M\int_a^b\min\{x-a,b-x\}\,dx.
\]
而
\[
\int_a^b\min\{x-a,b-x\}\,dx=\frac{(b-a)^2}{4}.
\]
所以
\[
M\ge\frac{4}{(b-a)^2}\left|\int_a^b f(x)\,dx\right|.
\]
取 \(\xi\) 使 \(|f'(\xi)|=M\)，便有
\[
|f'(\xi)|\ge
\frac{4}{(b-a)^2}\left|\int_a^b f(x)\,dx\right|
\ge
\frac1{(b-a)^2}\int_a^b f(x)\,dx.
\]
原命题成立。
\answer{结论成立，且可证明更强的绝对值估计。}
\examnote{端点为零的函数可用 \(\max|f'|\) 控制高度：\(|f(x)|\le M\min(x-a,b-x)\)。}
\end{solutionblock}

\begin{problemblock}
\textbf{例9.} 设函数 \(f(x)\) 有连续的二阶导数，
\[
f(1)=f(0)=1,\qquad
M=\max_{x\in[0,1]}|f''(x)|.
\]
\begin{enumerate}[label=(\arabic*), leftmargin=2.2em]
\item 证明：当 \(x\in[0,1]\) 时有 \(|f'(x)|\le\dfrac12M\)；
\item 证明：
\[
\left|\int_0^1f(x)\,dx\right|\le1+\frac{M}{8}.
\]
\end{enumerate}
\end{problemblock}
\begin{solutionblock}
\analysis{利用 \(f(1)=f(0)\) 得 \(\int_0^1f'(t)dt=0\)。第 (2) 问对 \(g=f-1\) 使用端点为零的高度估计。}
(1) 因
\[
\int_0^1f'(t)\,dt=f(1)-f(0)=0,
\]
所以
\[
f'(x)=\int_0^1\bigl[f'(x)-f'(t)\bigr]\,dt.
\]
又 \(|f''|\le M\)，由 Lagrange 中值定理得
\[
|f'(x)-f'(t)|\le M|x-t|.
\]
于是
\[
|f'(x)|
\le M\int_0^1|x-t|\,dt
=\frac{M}{2}\bigl[x^2+(1-x)^2\bigr]
\le\frac M2.
\]

(2) 令
\[
g(x)=f(x)-1.
\]
则 \(g(0)=g(1)=0\)，且由 (1)
\[
\max_{[0,1]}|g'(x)|=\max_{[0,1]}|f'(x)|\le\frac M2.
\]
因此
\[
|g(x)|\le\frac M2\min\{x,1-x\}.
\]
积分：
\[
\left|\int_0^1g(x)\,dx\right|
\le\frac M2\int_0^1\min\{x,1-x\}\,dx
=\frac M2\cdot\frac14=\frac M8.
\]
于是
\[
\left|\int_0^1f(x)\,dx\right|
=\left|1+\int_0^1g(x)\,dx\right|
\le1+\frac M8.
\]
\answer{两式均成立。}
\examnote{端点函数值相等时，\(f'\) 的平均值为零；再配合 \(|f''|\le M\) 可控制 \(f'\)。}
\end{solutionblock}

\begin{problemblock}
\textbf{例10.} 设 \(f\) 在 \([0,1]\) 有一阶导数且
\[
|f'(x)-f'(y)|\le |x-y|.
\]
证明：
\[
\left|\int_0^1f(x)\,dx-\frac{f(0)+f(1)}2\right|
\le \frac1{12}.
\]
\end{problemblock}
\begin{solutionblock}
\analysis{这是梯形公式误差估计。把误差写成 \(f'\) 与核函数的积分，再利用 \(f'\) 的 Lipschitz 条件。}
分部积分得
\[
\int_0^1 f(x)\,dx-\frac{f(0)+f(1)}2
=\int_0^1\left(\frac12-x\right)f'(x)\,dx.
\]
因为
\[
\int_0^1\left(\frac12-x\right)dx=0,
\]
所以可减去常数 \(f'(1/2)\)：
\[
\int_0^1\left(\frac12-x\right)f'(x)\,dx
=\int_0^1\left(\frac12-x\right)\left[f'(x)-f'\left(\frac12\right)\right]dx.
\]
由题设
\[
\left|f'(x)-f'\left(\frac12\right)\right|
\le\left|x-\frac12\right|.
\]
于是
\[
\begin{aligned}
\left|\int_0^1f(x)\,dx-\frac{f(0)+f(1)}2\right|
&\le
\int_0^1\left|x-\frac12\right|^2dx\\
&=\frac1{12}.
\end{aligned}
\]
\answer{不等式成立。}
\examnote{梯形误差核是 \(\frac12-x\)，其积分为零，所以可以减去任意常数，常取 \(f'(1/2)\)。}
\end{solutionblock}

\section{凸函数与 Lyapunov 型不等式}

\begin{problemblock}
\textbf{例11.} 设 \(f(x)\) 在 \([0,1]\) 上二阶可导，且满足
\[
f''(x)>0,\qquad \int_0^1f(x^2)\,dx=0,
\]
则（\quad）

A. \(f\left(\dfrac13\right)<0,\ 2f(0)+f(1)>0\)

B. \(f\left(\dfrac13\right)<0,\ f(0)+2f(1)>0\)

C. \(f\left(\dfrac13\right)<0,\ 2f(0)+f(1)>0\)

D. \(f\left(\dfrac13\right)<0,\ f(0)+2f(1)>0\)
\end{problemblock}
\begin{solutionblock}
\analysis{条件 \(f''>0\) 表示严格凸。第一个结论用 Jensen 不等式；第二个结论用凸函数位于端点弦下方。}
令 \(X\) 在 \([0,1]\) 上均匀分布，则
\[
\int_0^1 f(x^2)\,dx=E[f(X^2)]=0,
\qquad E[X^2]=\frac13.
\]
由于 \(f''>0\)，\(f\) 严格凸，由 Jensen 不等式
\[
f\left(\frac13\right)=f(E[X^2])<E[f(X^2)]=0.
\]
所以
\[
f\left(\frac13\right)<0.
\]

再由严格凸函数在端点弦下方，
\[
f(t)<(1-t)f(0)+tf(1)\qquad(0<t<1).
\]
取 \(t=x^2\) 并积分：
\[
0=\int_0^1f(x^2)\,dx
<
\int_0^1\bigl[(1-x^2)f(0)+x^2f(1)\bigr]dx.
\]
右边等于
\[
\frac23f(0)+\frac13f(1),
\]
故
\[
2f(0)+f(1)>0.
\]
\answer{A；题面中 C 与 A 为同一结论时也表示同一正确项。}
\examnote{严格凸时 Jensen 与弦不等式都取严格号，这是本题两个“大于/小于”严格成立的来源。}
\end{solutionblock}

\begin{problemblock}
\textbf{例12.} \(f(x)\) 在 \([0,1]\) 上具有二阶连续导数，且
\[
f(0)=f(1)=0,
\]
在 \((0,1)\) 内 \(f(x)\ne0\)。证明：
\[
\int_0^1\left|\frac{f''(x)}{f(x)}\right|\,dx\ge4.
\]
\end{problemblock}
\begin{solutionblock}
\analysis{这是二阶方程的 Lyapunov 型不等式。把 \(q(x)=-f''(x)/f(x)\)，用 Green 函数表示端点为零的函数。}
若积分为 \(+\infty\)，结论显然。下设该积分有限。

由于 \(f\) 在 \((0,1)\) 内不为零且连续，\(f\) 在 \((0,1)\) 内同号。令
\[
q(x)=-\frac{f''(x)}{f(x)}.
\]
则
\[
-f''(x)=q(x)f(x),\qquad f(0)=f(1)=0.
\]
端点为零的 Green 函数为
\[
G(x,t)=
\begin{cases}
t(1-x),&0\le t\le x\le1,\\
x(1-t),&0\le x\le t\le1.
\end{cases}
\]
于是
\[
f(x)=\int_0^1G(x,t)q(t)f(t)\,dt.
\]
取 \(x_0\in[0,1]\)，使
\[
|f(x_0)|=\max_{[0,1]}|f(x)|=:M_0>0.
\]
则
\[
M_0
\le\int_0^1G(x_0,t)|q(t)|\,|f(t)|\,dt
\le M_0\int_0^1G(x_0,t)|q(t)|\,dt.
\]
而对任意 \(x,t\in[0,1]\)，都有
\[
0\le G(x,t)\le\frac14.
\]
故
\[
1\le\frac14\int_0^1|q(t)|\,dt
=\frac14\int_0^1\left|\frac{f''(t)}{f(t)}\right|\,dt.
\]
即
\[
\int_0^1\left|\frac{f''(x)}{f(x)}\right|\,dx\ge4.
\]
\answer{不等式成立。}
\examnote{看到 \(f(0)=f(1)=0\) 且要估 \(\int|f''/f|\)，可联想到二阶边值问题的 Lyapunov 不等式。}
\end{solutionblock}

\begin{problemblock}
\textbf{例13.} \(f(x)\) 在 \((-\infty,+\infty)\) 上有二阶连续导数。证明：\(f''(x)\ge0\) 的充要条件是，对不同的实数 \(a,b\)，
\[
f\left(\frac{a+b}{2}\right)
\le
\frac1{b-a}\int_a^b f(x)\,dx.
\]
\end{problemblock}
\begin{solutionblock}
\analysis{必要性是凸函数的 Hermite-Hadamard 中点不等式；充分性对小区间作 Taylor 展开。}
先证必要性。若 \(f''(x)\ge0\)，则 \(f\) 为凸函数。记
\[
m=\frac{a+b}{2},\qquad h=\frac{b-a}{2}.
\]
对任意 \(t\in[0,h]\)，由凸性
\[
f(m)\le\frac{f(m-t)+f(m+t)}2.
\]
对 \(t\in[0,h]\) 积分：
\[
h f(m)\le\frac12\int_0^h[f(m-t)+f(m+t)]\,dt
=\frac12\int_{m-h}^{m+h}f(x)\,dx.
\]
由于 \(2h=b-a\)，得
\[
f\left(\frac{a+b}{2}\right)
\le
\frac1{b-a}\int_a^b f(x)\,dx.
\]

再证充分性。任取 \(x_0\in\mathbb{R}\)，对任意小 \(h>0\)，由题设
\[
f(x_0)\le \frac1{2h}\int_{x_0-h}^{x_0+h}f(x)\,dx.
\]
对 \(f\) 在 \(x_0\) 处作 Taylor 展开：
\[
f(x_0+u)=f(x_0)+f'(x_0)u+\frac12f''(x_0)u^2+o(u^2).
\]
在 \([-h,h]\) 上取平均，奇次项积分为零，得
\[
\frac1{2h}\int_{x_0-h}^{x_0+h}f(x)\,dx
=f(x_0)+\frac{f''(x_0)}6h^2+o(h^2).
\]
代入题设：
\[
0\le \frac{f''(x_0)}6h^2+o(h^2).
\]
令 \(h\to0^+\)，得
\[
f''(x_0)\ge0.
\]
因 \(x_0\) 任意，故 \(f''(x)\ge0\)。
\answer{充要条件成立。}
\examnote{“中点值不超过积分平均值”是凸函数的典型刻画；反向证明常用小区间 Taylor 平均。}
\end{solutionblock}

\begin{problemblock}
\textbf{例14.} 设当 \(x>-1\) 时，可微函数 \(f(x)\) 满足
\[
f'(x)+f(x)-\frac1{1+x}\int_0^x f(t)\,dt=0,
\qquad f(0)=1.
\]
试证明：当 \(x\ge0\) 时，有
\[
e^{-x}\le f(x)\le1.
\]
\end{problemblock}
\begin{solutionblock}
\analysis{设 \(F(x)=\int_0^x f(t)dt\)，把原方程中 \(f,F\) 的组合整理成一阶方程。}
令
\[
F(x)=\int_0^x f(t)\,dt.
\]
原方程等价于
\[
(1+x)f'(x)+(1+x)f(x)-F(x)=0.
\]
设
\[
H(x)=(1+x)f(x)-F(x).
\]
则
\[
H'(x)=f(x)+(1+x)f'(x)-F'(x)=(1+x)f'(x).
\]
由原方程又有
\[
(1+x)f'(x)=F(x)-(1+x)f(x)=-H(x).
\]
所以
\[
H'(x)=-H(x),\qquad H(0)=f(0)=1.
\]
于是
\[
H(x)=e^{-x}.
\]
即
\[
(1+x)f(x)-F(x)=e^{-x}.
\]
再由原方程
\[
f'(x)=\frac{F(x)}{1+x}-f(x)
=-\frac{e^{-x}}{1+x}<0\qquad(x\ge0).
\]
故 \(f\) 在 \([0,+\infty)\) 上单调递减，且 \(f(0)=1\)，所以
\[
f(x)\le1.
\]
另一方面，
\[
\bigl[f(x)-e^{-x}\bigr]'
=-\frac{e^{-x}}{1+x}+e^{-x}
=\frac{x e^{-x}}{1+x}\ge0.
\]
又 \(f(0)-e^0=0\)，因此
\[
f(x)-e^{-x}\ge0,\qquad x\ge0.
\]
即
\[
e^{-x}\le f(x)\le1.
\]
\answer{不等式成立。}
\examnote{含 \(\int_0^x f\) 的微分方程，常设 \(F=\int_0^x f\)，再寻找 \((1+x)f-F\) 这类可闭合的组合。}
\end{solutionblock}

\begin{problemblock}
\textbf{例15.} 设 \(f(x)\) 是 \([0,1]\) 上的可导函数，
\[
f(0)=f(1)=1,\qquad \max_{[0,1]}|f'(x)|=1,
\]
则（\quad）

A. \(\displaystyle \frac14<\int_0^1 f(x)\,dx<\frac12\)

B. \(\displaystyle \frac12<\int_0^1 f(x)\,dx<\frac34\)

C. \(\displaystyle \frac34<\int_0^1 f(x)\,dx<\frac54\)

D. \(\displaystyle \frac54<\int_0^1 f(x)\,dx<\frac74\)
\end{problemblock}
\begin{solutionblock}
\analysis{\(|f'|\le1\) 表示函数图像斜率受限。由两个端点同时控制，可得到上下包络。}
由 \(|f'|\le1\) 和 \(f(0)=1\)，得
\[
f(x)\ge1-x,\qquad f(x)\le1+x.
\]
由 \(|f'|\le1\) 和 \(f(1)=1\)，得
\[
f(x)\ge x,\qquad f(x)\le2-x.
\]
因此
\[
f(x)\ge\max\{1-x,x\},
\qquad
f(x)\le\min\{1+x,2-x\}.
\]
积分得
\[
\int_0^1\max\{1-x,x\}\,dx=\frac34,
\]
以及
\[
\int_0^1\min\{1+x,2-x\}\,dx=\frac54.
\]
所以
\[
\frac34\le\int_0^1f(x)\,dx\le\frac54.
\]
若左端等号成立，则必须
\[
f(x)=\max\{1-x,x\},
\]
该函数在 \(x=1/2\) 处不可导，与题设矛盾。右端等号同理要求
\[
f(x)=\min\{1+x,2-x\},
\]
也在 \(x=1/2\) 处不可导。因此两端均为严格不等式：
\[
\frac34<\int_0^1f(x)\,dx<\frac54.
\]
\answer{C}
\examnote{斜率受限题可从端点画“锥形上下包络”；等号常因包络在交点不可导而取不到。}
\end{solutionblock}
